\section{Literature}
\label{sec:literature}

\noindent Our paper contributes to the literature on the Great Depression by providing novel firm-level evidence about the causes of the delayed recovery. A pioneering contribution by \cite{Benmelech2018} shows that part of the employment decline from 1928 to 1933 was due to firms' inability to rollover maturing debt. We add to their contribution by identifying a complementary mechanism that contributed to the slow recovery of investment: the legal uncertainty introduced by the abrogation of gold clauses.

The legal uncertainty arising from the abrogation frames our study within the broader literature on economic uncertainty \citep{Bloom2009, Baker2016, Hassan2019}. We argue that leverage uncertainty created a debt overhang problem for firms with exposure to gold clauses. \cite{Hassan2019} measure firm-level exposure to policy shocks using textual analysis of earnings call transcripts. We adopt a similar firm-level approach, but identify exposure to legal uncertainty through the gold content of firms' liabilities. While \cite{Hassan2019} document the real effects of uncertainty in modern episodes such as Brexit, U.S. elections, and Italian constitutional reforms, we examine how uncertainty shaped firm decisions during the Great Depression, a pivotal historical episode. In related work, \cite{Bennett2025} extends the textual analysis methodology of \cite{Hassan2019} to measure firm-level litigation risk and its effects on corporate policies in contemporary settings.

Our work contributes to the empirical literature on debt overhang \citep{Myers1977} and firm outcomes. \cite{Lang1996} document a negative correlation between leverage and future hiring and investment. Subsequent studies aimed to establish causality by using exogenous shocks: \cite{Giroud2012} exploit debt restructuring of Austrian ski hotels to show that relief improves performance and reduces strategic defaults; \cite{Wittry2018} uses variation in mining reclamation liabilities across states; and \cite{Becker2012} exploit a legal change reducing equity–debt conflicts, finding effects consistent with \cite{Myers1977}. We contribute to this literature by showing that the debt overhang channel helps explain the delayed recovery of investment in the 1930s. The conceptual framework guiding our analysis builds on the dynamic overhang model of \cite{Hennessy2004}, who shows that defaultable long-term debt creates a wedge between marginal and average $Q$, which should be accounted for in conventional investment-$Q$ regressions.

The abrogation of gold clauses in sovereign debt contracts constituted a de facto US government default (\cite{Edwards2018}) and represents a major political intervention in private debt contracts (\cite{Bolton2002}). \cite{Edwards2015} show that, from 1933 to 1935, Treasury bond prices with and without gold clauses reflected uncertainty over the Supreme Court’s ruling, implying a significant perceived risk of unconstitutionality. \cite{Kroszner1999} finds that the Court's 1935 decision upholding abrogation led to broad increases in equity and corporate bond prices. The episode also informed theory: \citet{Bolton2002} model political intervention in private debt contracts, while \cite{Mian2014} use this precedent to study the welfare effects of debt forgiveness policies. We contribute to this literature by analyzing firms' investment and payout decisions during this episode using hand-collected firm-level data.